%%%%%%%%%%%%%%%%%%%%%%%%%%%%%%%%%
%
%       EXERCÍCIO
%
%%%%%%%%%%%%%%%%%%%%%%%%%%%%%%%%%

\definevalorquestao

\begin{exercicioBanco}[\valorquestao]
    Para comparar as médias de duas populações Normais, amostras aleatórias foram obtidas. Sabe-se que as variâncias populacionais são diferentes, sendo seus valores desconhecidos.

    \begin{center}
    \begin{tabular}{|c|cccccccc|}
        \hline
        & \multicolumn{8}{c|}{Observações} \\
        \hline
        Amostra I  & 7 & 9 & 3 & 8 & 11 & 5 & 9 & {} \\
        Amostra II & 2 & 7 & 5 & 15 & 9 & 16 & 8 & {} \\
        \hline
    \end{tabular}
    \end{center}

    O que pode ser dito a respeito das médias das populações, com
    \(\alpha = 0{,}05\)?
\end{exercicioBanco}

